\begin{figure}[t]
\centering
\begin{tikzpicture}
  \begin{groupplot}[
      group style={group size=3 by 1, horizontal sep=10mm},
      width=3.5cm, height=3.5cm, scale only axis,
      unbounded coords=jump,
      title style={font=\footnotesize, yshift=2pt}]
  \nextgroupplot[
      axis line style={cgrey, line width=0.4pt},
      tick style={cgrey, line width=0.4pt},
      label style={font=\footnotesize}, tick label style={font=\scriptsize},
      xlabel={index $i$}, ylabel={$\lambda_i(L_{\mathrm{norm}})$},
      xmin=0.5, xmax=10.5, ymin=-0.004, ymax=0.09,
      xtick={2,4,6,8,10}, scaled y ticks=false,
      yticklabel style={/pgf/number format/.cd, fixed, precision=2},
      title={(a) eigengap heuristic}]
    \addplot[cblue, mark=*, mark size=1.1pt, line width=0.7pt]
        table[x=i, y=lambda, col sep=comma] {\dataroot/spectrum_normalized.csv};
    \draw[cred, line width=1.4pt] (axis cs:4,0.0321) -- (axis cs:5,0.0449);
    \node[font=\scriptsize, cred, anchor=north west] at (axis cs:4.5,0.030)
        {$k^\star=4$};
  \nextgroupplot[hide axis, xmin=-1.12, xmax=1.12, ymin=-1.12, ymax=1.12,
                 enlargelimits=false, title={(b) projection on $\bm v_2,\bm v_3$}]
    \addplot[cluster marks tiny] table[x=x, y=y, meta=class, col sep=comma]
        {\dataroot/cluster_embedding.csv};
  \nextgroupplot[hide axis, xmin=-1.12, xmax=1.12, ymin=-1.12, ymax=1.12,
                 enlargelimits=false, title={(c) the clusters in space}]
    \addplot[cluster marks tiny] table[x=x, y=y, meta=class, col sep=comma]
        {\dataroot/cluster_structure.csv};
  \end{groupplot}
\end{tikzpicture}
\caption{Spectral clustering of the haemoglobin contact network of
Figure~\ref{fig:examples}(c). (a)~the largest gap in the spectrum of
$L_{\mathrm{norm}}$, among the candidate gaps after $k=2,\dots,9$, occurs
after the fourth eigenvalue, so \eqref{eq:eigengap} returns $k^\star=4$
within the prespecified coarse-partition range;
(b)~the $574$ residues projected on the first two nontrivial generalized eigenvectors $\bm v_2,\bm v_3$ of \eqref{eq:koren}, with a marker per cluster;
(c)~the same clusters at the residue positions in space, projected on their
two principal axes. Compare with Figure~\ref{fig:examples}(c), which is the
same projection coloured by the true chain.}
\label{fig:clustering}
\end{figure}
